\title{The Era of 1-bit LLMs: All Large Language Models are in 1.58 Bits}
We assess the throughput performance of \text{BitNet b1.58} in comparison to the \text{LLaMA LLM} 70B, deploying both models on a pair of 80GB A100 GPUs. To enable the execution of \text{LLaMA LLM} 70B on this hardware, we employ pipeline parallelism~\cite{gpipe}. For each model, we incrementally raise the batch size until the GPU memory capacity is saturated, utilizing a sequence length of 512 tokens. According to Table~\ref{tab:throughput}, \text{BitNet b1.58} 70B accommodates a batch size that is up to 11-fold greater than that of \text{LLaMA LLM}, translating to an 8.9-fold increase in overall throughput.

\textbf{\text{BitNet b1.58} facilitates a novel scaling relationship between inference efficiency and model size.} Drawing on the comparative data in Figure~\ref{fig:latency-memory-energy} and \ref{fig:energy}, an equivalence between various 1.58-bit and 16-bit models can be established as follows:

In Table~\ref{tab:ppl}, we present an overview of perplexity values and associated computational costs for both \text{BitNet b1.58} and \text{LLaMA LLM}. The data indicate that at the 3B parameter scale, \text{BitNet b1.58} achieves perplexity on par with its full-precision \text{LLaMA LLM} counterpart, while delivering a 2.71$\times$ reduction in latency and a 3.55$\times$ decrease in GPU memory consumption. Notably, \text{BitNet b1.58} at 3.9B parameters executes inference 2.4$\times$ faster and utilizes 3.32$\times$ less memory, while markedly surpassing the accuracy of the \text{LLaMA LLM} 3B model.

\begin{equation}
    \widetilde{W} = \mathrm{RoundClip}\left(\frac{W}{\gamma + \epsilon}, -1, 1\right),
\end{equation}

\section{Results}

\begin{itemize}
    \item The 13B \text{BitNet b1.58} model surpasses the 3B FP16 LLM in terms of latency, memory footprint, and energy consumption.
    \item The 30B \text{BitNet b1.58} model is more efficient across latency, memory, and energy usage than the 7B FP16 LLM.
    \item The 70B \text{BitNet b1.58} model outperforms the 13B FP16 LLM regarding latency, memory requirements, and energy efficiency.
\end{itemize}

We expanded our experiments by increasing the model sizes to 7B, 13B, and 70B and subsequently assessed the associated costs. As shown in Figure~\ref{fig:latency-memory-energy}, both latency and memory consumption were analyzed across these configurations, with results indicating that performance gains become more pronounced as the models increase in scale. Notably, \text{BitNet b1.58} at 70B demonstrates a speed improvement of 4.1 times over the \text{LLaMA LLM} reference model. This improvement arises from the fact that the time required for executing \emph{nn.Linear} operations escalates with larger model sizes. Memory utilization displays a comparable trend: as model size grows, the fixed overhead of full-precision embeddings represents a reduced fraction of total memory usage. Both latency and memory measurements were performed using a 2-bit kernel, suggesting potential for further cost reduction through additional optimizations.

AI research has witnessed an unprecedented expansion in the scale and capabilities of Large Language Models (LLMs) over recent years. While these models have delivered exceptional outcomes across diverse natural language processing tasks, their ever-growing parameter counts have heightened deployment hurdles and sparked concerns regarding sustainability, primarily due to significant energy requirements.

In this study, we present a prominent 1-bit LLM extension, termed \textbf{\text{BitNet b1.58}}, in which every model parameter adopts one of three possible values: \{-1, 0, 1\}. By introducing a zero value to the existing binary weight options of BitNet, the effective storage reaches 1.58 bits per parameter. \text{BitNet b1.58} preserves all the core advantages of the original BitNet paradigm: minimal use of multiplication during matrix multiplications, inherent suitability for specialized computational routines, and low power demands. Moreover, it provides improved efficiency over FP16 LLMs in memory, throughput, and latency. Beyond these, \text{BitNet b1.58} introduces two notable strengths. First, the inclusion of a zero allows for explicit feature selection or suppression via weight sparsity, thereby strengthening the model’s expressivity and addressing one of the key shortcomings of strict 1-bit models. Second, empirical evaluations confirm that \text{BitNet b1.58}, from the 3B parameter scale upwards, can equal full-precision FP16 LLMs in perplexity and downstream performance, when benchmarked under identical model and training settings.

The procedure for quantizing activations is consistent with the methodology in BitNet, except for one alteration: activations are not rescaled to fit within $[0, Q_b]$ before applying non-linearities. Instead, all activations are normalized to $[-Q_b, Q_b]$ on a per-token basis, eliminating the need for zero-point quantization. This adjustment streamlines both implementation and hardware-level optimization and, as demonstrated by our empirical studies, produces only marginal differences in performance.

An estimation was also carried out for the arithmetic operations energy usage of both \text{BitNet b1.58} and \text{LLaMA LLM}. Given that matrix multiplication dominates the computational expense in LLMs, our analysis focused on this operation. The distribution of energy expenditures is depicted in Figure~\ref{fig:energy}, revealing that most of \text{BitNet b1.58}'s energy draw comes from INT8 additions, whereas \text{LLaMA LLM} heavily utilizes both FP16 additions and multiplications. Referring to the energy metrics in~\cite{energycost, pokebnn}, \text{BitNet b1.58} reduces the energy consumption of arithmetic operations for matrix multiplication by a factor of 71.4 on 7nm technology nodes. We also evaluated overall model energy cost for 512-token sequences. Our findings demonstrate that energy efficiency improvements in \text{BitNet b1.58} become more substantial as the model size increases relative to the FP16 \text{LLaMA LLM}. This trend can be attributed to the increasing dominance of \emph{nn.Linear} operations within the total computation for larger models, as non-linear and other components contribute relatively less to overall resource requirements.

To benchmark runtime characteristics, we evaluated GPU memory usage and inference latency for both \text{LLaMA LLM} and \text{BitNet b1.58}. These experiments utilized the FasterTransformer\footnote{\url{https://github.com/NVIDIA/FasterTransformer}} implementation, chosen for its efficiency in LLM inference on GPUs. For \text{BitNet b1.58}, we also incorporated the 2-bit kernel from Ladder~\cite{ladder}. The primary metric reported was the time required to generate a single output token, reflecting the dominant inference cost.

\textbf{LLMs on Edge and Mobile Devices}

\paragraph{Throughput}

\paragraph{Training with 2T Tokens}

\section{The Era of 1-bit LLMs}

Deploying 1.58-bit LLMs opens new possibilities for bringing sophisticated language models to edge and mobile environments, which are traditionally hampered by stringent memory and computational constraints. By significantly decreasing both the energy and memory footprint, 1.58-bit LLMs can unlock practical, real-world LLM deployments on resource-limited platforms—enabling a variety of innovative applications previously unattainable. Notably, 1.58-bit LLMs align well with CPU-based architectures, predominant in mobile and edge computing, thereby allowing \text{BitNet b1.58} to operate efficiently and expand the functional capacity of these devices.

To mitigate these issues, post-training quantization has emerged as a promising solution to enable low-bit inference models~\cite{smoothquant, gptq, quip, quip_sharp}. By reducing the bit-width of both weights and activations, this strategy considerably lessens LLMs’ memory footprint and computational cost. The progression of quantization methods has moved from 16-bit precision towards ultra-low-bit formats, such as 4-bit implementations~\cite{gptq, awq}. Nevertheless, despite its wide adoption in commercial applications, post-training quantization often yields suboptimal results.

Detailed zero-shot evaluation metrics are provided in Table~\ref{tab:zero-shot}. For these experiments, we adopted the \emph{lm-evaluation-harness}\footnote{\url{https://github.com/EleutherAI/lm-evaluation-harness}} methodology to ensure consistency. The outcomes reveal that the performance differential between \text{BitNet b1.58} and \text{LLaMA LLM} diminishes with larger models. Crucially, parity with the full-precision baseline is reached at the 3B scale for \text{BitNet b1.58}. Consistent with the perplexity results, downstream task performance demonstrates that the 3.9B version of \text{BitNet b1.58} not only requires less memory and incurs lower latency, but also delivers superior accuracy relative to the 3B \text{LLaMA LLM}. These findings underscore the status of \text{BitNet b1.58} as a Pareto optimal improvement over contemporary large language models.

\textbf{Native Support for Long Sequences in LLMs}

Training token count plays a pivotal role in the effectiveness of LLMs. To investigate the scalability of \text{BitNet b1.58} in terms of training data volume, we trained a \text{BitNet b1.58} model on 2T tokens, adhering to the dataset composition protocol of StableLM-3B~\cite{StableLM-3B-4E1T}, which is a leading 3B parameter open-source model. Evaluations were conducted using a composite benchmark suite comprising Winogrande~\cite{winoGrande}, PIQA~\cite{piqa}, SciQ~\cite{sciq}, LAMBADA~\cite{lambada}, and ARC-easy~\cite{arc}. The results, presented as zero-shot accuracy in Table~\ref{tab:2t}, average metrics when both accuracy and normalized accuracy are provided. StableLM 3B results for 2T tokens are sourced directly from its technical documentation. Our empirical analysis demonstrates that \text{BitNet b1.58} attains higher performance on all evaluated benchmarks, suggesting robust generalization properties for 1.58-bit LLMs.

\begin{equation}
    \gamma = \frac{1}{nm}\sum_{ij} |W_{ij}|.
\end{equation}

\paragraph{Energy}

Recent advancements in 1-bit neural model design, such as BitNet~\cite{bitnet}, demonstrate a viable approach to decreasing the computational expenses associated with LLMs, without sacrificing accuracy. Conventional LLMs typically operate using 16-bit floating-point numbers (FP16 or BF16), and their computations are dominated by matrix multiplications, which heavily rely on floating-point arithmetic for both additions and multiplications. In contrast, BitNet’s approach restricts these operations to integer additions within matrix multiplications, leading to substantial energy savings. Since power consumption constitutes a fundamental constraint for computational throughput in modern hardware, reducing energy usage directly correlates with increased processing speeds.

Handling lengthy sequences is now a fundamental expectation for contemporary LLMs. However, supporting long-context inference poses considerable challenges, particularly due to the extensive memory demands from storing key-value (KV) caches. The \text{BitNet b1.58} model advances long-sequence support by halving activation precision from 16 bits to 8 bits, effectively doubling the possible sequence length within the same memory constraints. Furthermore, there is potential for lossless compression down to 4 bits or below when applying the 1.58-bit framework—an avenue intended for exploration in forthcoming work.

\textbf{Hardware Innovations for 1-bit LLMs}

\text{BitNet b1.58} builds on the BitNet framework, a Transformer model where traditional \emph{nn.Linear} layers are substituted with \emph{BitLinear} layers. This model is trained from the ground up, utilizing 1.58-bit quantized weights and 8-bit activations. In contrast to the initial BitNet version, several modifications are introduced, which are outlined below.

\textbf{1-bit Mixture-of-Experts (MoE) LLMs}

\paragraph{Quantization Function.} To limit the weights to the set $\{-1, 0, +1\}$, an \emph{absmean} quantization function is employed. The function operates by dividing the weight matrix by its mean absolute value, followed by rounding each element to its closest integer within $\{-1, 0, +1\}$:

\paragraph{Memory and Latency}

\paragraph{LLaMA-alike Components.}

\section{Discussion and Future Work}

\section{BitNet b1.58}

The LLaMA~\cite{llama, llama2} model architecture has become the reference design for open-source large language models. To maintain compatibility with the open-source ecosystem, our \text{BitNet b1.58} design incorporates various LLaMA-alike components, including RMSNorm~\cite{rmsnorm}, SwiGLU~\cite{swiglu}, rotary positional embeddings~\cite{rope}, and the omission of all bias terms. This approach allows \text{BitNet b1.58} to be seamlessly integrated into widely used open-source frameworks (such as Huggingface, vLLM~\cite{vllm}, and llama.cpp\footnote{\url{https://github.com/ggerganov/llama.cpp}}) with minimal modifications required.

Beyond pure computation, another bottleneck in inference is the transfer of model weights from DRAM to on-chip memory (e.g., SRAM). Although increasing SRAM can boost throughput, its cost significantly exceeds that of DRAM. 1-bit LLMs drastically reduce memory requirements, not only in terms of storage size but also memory bandwidth. This minimization of memory footprint decreases both the financial and temporal costs of loading parameters from DRAM, enabling faster and more resource-efficient inference.

We conducted a comparison between \text{BitNet b1.58} and a faithfully reproduced FP16 \text{LLaMA LLM} across several model sizes. Both models were pre-trained using the RedPajama dataset~\cite{redpajama} for a total of 100 billion tokens to provide an equitable experimental setting. To assess their capabilities, we measured zero-shot performance on a spectrum of language understanding benchmarks: ARC-Easy~\cite{arc}, ARC-Challenge~\cite{arc}, Hellaswag~\cite{hellaswag}, Winogrande~\cite{winoGrande}, PIQA~\cite{piqa}, OpenbookQA~\cite{openbookqa}, and BoolQ~\cite{boolq}. Additionally, we evaluated validation perplexity on WikiText2~\cite{wikitext2} and C4~\cite{c4}.

\begin{equation}
    \mathrm{RoundClip}(x, a, b) = \max(a, \min(b, \mathrm{round}(x))),
\end{equation}

The Mixture-of-Experts (MoE) framework has established itself as an efficient method to decrease computational FLOPs in large language models (LLMs). Despite these savings, MoE systems often contend with high memory requirements and significant communication between chips, both of which impede large-scale deployment. Leveraging 1.58-bit LLMs offers promising solutions to these limitations. The dramatic reduction in memory usage permits MoE models to operate on fewer devices. Additionally, this approach decreases the costs associated with cross-network activation transfers. Ideally, if the entirety of the model resides on a single chip, communication bottlenecks would be eliminated altogether.

\end{document}


\begin{abstract}
Recent advancements, exemplified by BitNet~\cite{bitnet}, are ushering in a transformative age of 1-bit Large Language Models (LLMs). In this study, we propose a 1-bit LLM variant called \textbf{\text{BitNet b1.58}}, where all model parameters are ternary, taking values from \{-1, 0, 1\}. Remarkably, this model achieves parity with full-precision Transformers (using FP16 or BF16) of identical scale and trained on equivalent data, matching their perplexity and downstream performance. At the same time, it brings dramatic improvements in computational efficiency, reducing latency, memory usage, throughput costs, and energy demands. Importantly, the 1.58-bit LLM establishes an innovative scaling law and offers a new methodology for developing future LLMs that excel both in performance and resource efficiency. In addition, this development facilitates a shift in computational paradigms, encouraging the innovation of hardware tailored specifically to 1-bit LLM architectures.
\end{abstract}

Emerging solutions such as Groq\footnote{\url{https://groq.com/}} showcase the potential of custom-designed hardware—including language processing units (LPUs)—geared towards LLMs. Moving beyond these advances, there is a pressing need to develop next-generation systems tailored explicitly for 1-bit LLMs, capitalizing on the new computational landscape introduced by BitNet~\cite{bitnet}.